\chapter{Cadre Théorique du Système d'Information}

\section{Introduction}
Avant d'exposer le contexte organisationnel de Naftal et la solution GIA, il convient de
rappeler les notions qui fondent tout projet de système d'information (SI) en entreprise.
Ce chapitre prolonge les développements classiques du management des SI — tels qu'on les
retrouve notamment dans les travaux sur la fonction support et la DCSI au sein de
Naftal~\cite{meghari2014,ssiouani2010} — et les transpose au périmètre d'un projet
d'ingénierie logicielle (helpdesk applicatif). Nous y précisons : la nature de l'information, les
définitions du SI, ses fonctions et composantes, la distinction SI\,/\,informatique, les rôles
stratégiques, l'organisation de la fonction SI, puis des \textbf{repères} sur le modèle
d'évaluation fonctionnelle (MEF)~\cite{autissier2008}, utiles pour comprendre pourquoi des
outils de traçabilité et de pilotage du support s'inscrivent dans une logique de performance des
fonctions transverses.

\section{De la donnée à l'information, du système au système d'information}

\subsection{Données et information}
Les auteurs classiques du management des systèmes d'information distinguent habituellement
la \textbf{donnée}, à l'état brut, de l'\textbf{information}, c'est-à-dire une donnée contextualisée,
structurée et rendue exploitable pour un décideur ou un opérateur~\cite{reix1995}. Les données
relèvent d'événements ou de mesures ; l'information naît lorsqu'elles sont organisées de sorte
que les acteurs puissent les comprendre et agir. Le SI assure précisément cette transformation :
collecte, mise en forme, stockage et diffusion. Pour l'organisation, la qualité de l'information
(pertinence, fiabilité, disponibilité, coût maîtrisé) conditionne la qualité des décisions et la
réactivité face aux incidents ou aux changements.

\subsection{Le système au sens organisationnel}
Le traitement des flux d'information internes et externes suppose des composantes reliées entre
elles et orientées vers un objectif commun : on parle alors de \textbf{système}, avec des entrées,
un traitement et des sorties. On peut modéliser l'entreprise comme un système ouvert : un
\textbf{système opérant} produit les biens ou services, un \textbf{système de pilotage} fixe les
objectifs et alloue les moyens, et le \textbf{système d'information} assure la circulation de
l'information entre ces niveaux et avec l'environnement. Le pilotage transmet des directives à
l'opérant via le SI ; l'opérant rend compte des réalisations via le même relais. Une défaillance
du support applicatif ou réseau n'est donc pas un simple « problème technique » : elle rompt la
chaîne d'information sur laquelle reposent coordination et contrôle.

Dans le cadre du projet GIA, cette logique systémique se traduit par la nécessité de fiabiliser les
flux d'information entre demandeurs, techniciens et administrateurs, afin de préserver la
continuité opérationnelle.

\section{Panorama des définitions du système d'information}
La plupart des définitions mettent l'accent soit sur les \textbf{composantes} (matériel, logiciel,
personnel, données, procédures), soit sur les \textbf{fonctions} (acquérir, traiter, stocker,
communiquer). Le tableau~\ref{tab:defsi} en propose une synthèse indicative, sans exhaustivité
bibliographique, afin de montrer la diversité des angles d'analyse retenus dans la littérature
francophone (formulations adaptées et condensées).

\begin{table}[H]
\centering
\footnotesize
\begin{tabularx}{\textwidth}{@{}lX@{}}
\toprule
\textbf{Auteurs (indicatif)} & \textbf{Propos central sur le SI} \\
\midrule
Vidal, Planeix et al. &
Ensemble organisé de ressources (matériels, logiciels, personnels, données, procédures)
permettant d'acquérir, de traiter, de stocker et de communiquer des informations sous
formes variées (données, textes, images, sons) dans des organisations. \\
\midrule
Reix et Rowe (2002) &
Ensemble d'acteurs sociaux qui mémorisent et transforment des représentations via des
technologies de l'information et des modes opératoires. \\
\midrule
Laudon et Laudon &
Ensemble de moyens humains, matériels et organisationnels permettant de recueillir, de
traiter, de stocker et de transmettre l'information au sein d'une organisation. \\
\midrule
Kalika et collaborateurs &
Système intégré utilisateur--machine produisant de l'information pour assister les
fonctions d'exécution, de gestion et de décision, à l'aide d'équipements, logiciels, bases
de données, procédures manuelles et modèles d'analyse. \\
\midrule
Ngang Bilonga, Jacques &
Ensemble structuré de données et de moyens techniques pour générer, mémoriser, traiter,
transférer et exploiter des informations dans un cadre d'objectifs définis. \\
\midrule
Le Moigne &
Fonction consistant à produire et enregistrer les informations représentatives de
l'activité de l'opération, puis à les mettre à disposition, de manière interactive lorsque
possible, du système de décision. \\
\bottomrule
\end{tabularx}
\caption{Illustration de définitions du système d'information (synthèse documentaire)}
\label{tab:defsi}
\end{table}

À partir de ces approches, nous retenons une définition opérationnelle alignée sur Reix~\cite{reix1995}
déjà citée en introduction : le SI est un ensemble hétérogène d'éléments humains et techniques,
en interaction, qui transforme des entrées en sorties informationnelles au service de la gestion.
Il se caractérise par son hétérogénéité (nature et forme de l'information), sa complexité (flux
internes et externes) et son caractère évolutif (adaptation au contexte).

Appliquée à GIA, cette définition rappelle que la plateforme ne se limite pas à un logiciel : elle
structure une organisation du support, des procédures et des responsabilités autour de
l'information d'incident.

\section{Fonctions et composantes du système d'information}

\subsection{Les quatre fonctions de base}
On distingue classiquement quatre fonctions, présentes à l'échelle de la plateforme GIA :
\begin{itemize}
    \item \textbf{Acquisition :} collecte et saisie — formulaire de ticket, pièces jointes,
    identification de l'auteur.
    \item \textbf{Traitement :} application de règles — priorités, statuts, contrôle d'accès,
    désassignation lors d'un retour en file d'attente.
    \item \textbf{Stockage :} persistance — base SQL Server, journal d'audit.
    \item \textbf{Communication / diffusion :} mise à disposition via l'intranet — listes,
    tableaux de bord, historique chronologique.
\end{itemize}
L'objectif est de préserver, exploiter et échanger l'information pour automatiser des tâches
répétables ou fournir aux utilisateurs les éléments nécessaires à une action rapide et éclairée.

La plateforme GIA incarne précisément ces quatre fonctions : elle acquiert les données via le
formulaire de ticket, les traite selon les règles métier (priorité, assignation, statut), les
stocke dans SQL Server et les diffuse via l'intranet aux acteurs concernés.

\subsection{Les composantes}
Un SI mobilise des \textbf{ressources humaines} : utilisateurs et cadres qui consomment ou
produisent l'information, et spécialistes (concepteurs, analystes, exploitants) qui conçoivent
et font fonctionner le dispositif. Les \textbf{moyens matériels} couvrent postes, serveurs,
réseaux et supports d'information. Les \textbf{moyens logiciels} regroupent programmes
d'application et d'infrastructure ainsi que les \textbf{procédures} (saisie, contrôles,
corrections). Les \textbf{données} constituent la matière première structurée (identités,
tickets, journaux). Le SI ne résulte pas de la juxtaposition de ces briques : il suppose une
construction visant explicitement les objectifs des utilisateurs futurs.

\section{Système d'information, informatique, maîtrise d'ouvrage et maîtrise d'œuvre}
L'\textbf{informatique} désigne souvent l'outil, le centre de coûts techniques, la maîtrise des
langages et du paramétrage. Le \textbf{système d'information} renvoie à une vision plus globale :
philosophie de gestion, chaîne de valeur informationnelle, processus et acteurs. La relation est
fréquemment qualifiée de \emph{demande} / \emph{offre} ou de client / fournisseur. La
\textbf{maîtrise d'ouvrage (MOA)} est responsable de l'efficacité organisationnelle et des méthodes
de travail autour du SI ; elle s'appuie sur une \textbf{maîtrise d'œuvre (MŒ)} interne ou externe
pour obtenir matériels, logiciels et services. Les informaticiens relèvent typiquement de la MŒ ;
la définition des besoins métier et la valeur créée relèvent davantage de la MOA et des directions
opérationnelles. Un helpdesk applicatif comme GIA matérialise cette articulation : les métiers
expriment le besoin de traçabilité ; la réalisation technique en découle.

\section{Déterminants de l'intensité d'usage des systèmes d'information}
La littérature retient souvent trois familles de facteurs : (1) la \textbf{taille} de l'entreprise —
plus elle est grande, plus le recours au SI est fréquent ; (2) le caractère \textbf{immatériel} de
l'activité — davantage d'information à traiter ; (3) le degré d'\textbf{internationalisation} —
SI plus élaboré pour coordonner des périmètres étendus. Naftal combine grande taille nationale,
activités logistiques et commerciales fortement informatisées et enjeux de groupe : l'intensité
d'usage du SI et les attentes en matière de support sont élevées.

\section{Typologies des systèmes d'information}
Les SI se classent selon plusieurs critères (niveau de décision, fonction métier, finalité).
Le tableau~\ref{tab:typosi} en donne une grille classique, utile pour situer GIA : il s'agit
principalement d'un système d'appui à l'\textbf{exploitation} et à la \textbf{gestion des anomalies}
au sein de la fonction maintenance applicative / relation utilisateurs, avec des effets sur le
pilotage (KPIs).

\begingroup
\setlength{\tabcolsep}{5pt}
\renewcommand{\arraystretch}{1.15}
\small
\begin{longtable}{@{}>{\raggedright\arraybackslash}p{0.34\textwidth}>{\raggedright\arraybackslash}p{0.62\textwidth}@{}}
\caption{Typologies usuelles des systèmes d'information (adaptation documentaire)\label{tab:typosi}} \\
\toprule
\textbf{Critère / type} & \textbf{Rôle attendu} \\
\midrule
\endfirsthead
\multicolumn{2}{c}{\tablename\ \thetable{} — \textit{suite}} \\
\toprule
\textbf{Critère / type} & \textbf{Rôle attendu} \\
\midrule
\endhead
\midrule
\multicolumn{2}{r}{\textit{Suite page suivante}} \\
\endfoot
\bottomrule
\endlastfoot
Systèmes du niveau \textbf{stratégique} &
Aider les cadres supérieurs à traiter des questions de long terme et à réagir aux tendances
de l'entreprise et de son environnement. \\
Systèmes du niveau \textbf{gestion} &
Soutenir les cadres intermédiaires dans le suivi, le contrôle, la décision non routinière et
l'administration. \\
Systèmes du niveau \textbf{opérations} &
Aider les cadres opérationnels à suivre activités et transactions élémentaires (ventes, paie,
etc.). \\
Système d'information \textbf{marketing} &
Produire un flux d'information structuré, interne et externe, au service des décisions
marketing. \\
Système d'information de \textbf{production} &
Organiser les activités de production pour satisfaire les commandes. \\
Système d'information \textbf{comptable / financier} &
Formaliser la gouvernance financière et la traçabilité des flux monétaires. \\
Système d'information des \textbf{RH} &
Couvrir le domaine des ressources humaines (paie, compétences, formation, etc.). \\
Système d'information pour \textbf{dirigeants} &
Fournir une vision synthétique et actualisée des activités et opérations. \\
Système d'information de \textbf{gestion du SI} &
Identifier les activités liées au pilotage du SI lui-même et aux métiers associés. \\
Système d'information d'\textbf{aide à la décision} &
Suivi opérationnel, conservation de l'historique, production de statistiques et de rapports. \\
Système de \textbf{traitement des transactions} &
Garantir des propriétés informatiques fortes (cohérence, intégrité) sur les données
transactionnelles. \\
\end{longtable}
\endgroup

\section{Rôles stratégiques du système d'information}
Au-delà de l'automatisation, le SI joue au moins trois rôles structurants.

\subsection{Instrument de partage et de connexion}
L'efficacité de la décision et la rapidité de réaction dépendent de la qualité du partage :
rapidité de transmission, fiabilité (limitation des « bruits » et déformations), complétude
(absence d'omissions qui biaiseraient la décision) et adéquation au besoin du destinataire selon
sa fonction.

\subsection{Mémoire de l'organisation}
Une organisation qui perd sa mémoire perd une partie de son savoir et de son savoir-faire.
Capitaliser les informations — et les restituer au bon moment — est une fonction explicitement
assumée par les bases de données et, dans notre cas, par le journal \texttt{incident\_logs}.

La table \texttt{incident\_logs} de GIA joue exactement ce rôle de mémoire organisationnelle :
chaque action est horodatée et attribuée, constituant un historique permanent opposable lors d'un
audit.

\subsection{Instrument de mise en forme}
Le SI doit fournir l'information sous une forme compréhensible et adaptée au niveau
hiérarchique : tableaux de bord pour l'administrateur, vues filtrées pour le technicien, suivi
personnel pour le déclarant.

\section{Acteurs internes de la fonction système d'information}
Les entreprises mobilisent une diversité de métiers : \textbf{programmeurs} (instructions
logicielles), \textbf{analystes} (traduction des besoins métiers en exigences d'information),
\textbf{gestionnaires de SI} (encadrement des équipes, projets, infrastructures,
télécommunications, bases de données), \textbf{superviseurs} des opérations et de la saisie,
\textbf{spécialistes externes} (constructeurs, éditeurs, consultants). Les \textbf{utilisateurs
finaux}, hors service SI, participent de plus en plus à la conception. Dans les grandes
structures, un \textbf{responsable du SI} de haut niveau peut piloter l'usage stratégique des
technologies ; des rôles spécialisés (sécurité, connaissances, protection des données) complètent
souvent l'organigramme. Le service SI agit aussi comme \textbf{agent de changement} : nouvelles
stratégies, produits et services fondés sur l'information.

\section{Formes d'organisation de la fonction SI}
Trois schémas reviennent fréquemment dans la littérature. (1) Structure \textbf{décentralisée} :
chaque grande fonction (marketing, production, etc.) possède son propre service SI — réactivité
métier mais risque d'incompatibilités technologiques et de surcoûts. (2) Service SI
\textbf{centralisé} : une direction unique arbitre les choix pour l'ensemble de l'entreprise —
meilleure cohérence et plans de développement à long terme. (3) Modèle \textbf{hybride} :
divisions autonomes mais soumises à un noyau central qui fixe normes, achats et orientations
d'architecture. Naftal illustre une combinaison de \textbf{pilotage central} (DCSI) et de
\textbf{réalités de terrain} par branches.

\section{Missions usuelles de la direction des systèmes d'information}
La direction des SI interagit avec la direction générale et les métiers. Elle assure en pratique
la coordination transversale, le contrôle des acteurs internes et externes, la proposition
d'évolutions alignées sur l'organisation, la cohérence stratégique et la rationalisation
technico-budgétaire. Sur le plan organisationnel, on retrouve souvent une \textbf{cellule
études / projets} (applications, maintenance applicative), une \textbf{cellule exploitation}
(mise à disposition des services, support bureautique et applicatif) et une \textbf{cellule
infrastructure et télécommunications} (réseaux, messagerie). La centralisation ou la
décentralisation de ces cellules dépend de la taille et du modèle du groupe.

\section{Profil du directeur des systèmes d'information}
Le directeur des SI doit combiner compétences technologiques, communicationnelles et
manageriales : fixer les orientations fonctionnelles et la stratégie SI, superviser conception,
mise en œuvre et maintien en conditions opérationnelles (qualité, sécurité, coûts, délais),
piloter des projets, gérer les budgets. L'activité s'exerce en lien avec prestataires, MOA, MŒ
et direction générale, avec des variations selon la structure.

\section{Le modèle d'évaluation fonctionnelle (MEF) : repères}
La fonction SI est une \textbf{fonction support} : elle ne constitue pas un centre de profit mais
contribue à la valeur par les services rendus aux métiers. Évaluer sa performance ne se réduit pas
à des ratios financiers ; le \textbf{modèle d'évaluation fonctionnelle (MEF)}~\cite{autissier2008}
propose une lecture en \textbf{quatre pôles} — activités, compétences, organisation, satisfaction
des clients internes — permettant des diagnostics comparés à un référentiel et l'identification de
pistes d'amélioration (questionnaires et baromètres dans les études de terrain).

Le pôle \textbf{activités} s'appuie sur un référentiel décrivant ce que la fonction SI doit
accomplir. Dans la version appliquée à la fonction SI, on distingue notamment cinq
\textbf{rubriques métiers}~\cite{autissier2008,meghari2014} :
\begin{enumerate}
    \item \textbf{Pilotage du département SI} — stratégie, schéma directeur, veille, tableaux de
    bord de la fonction, gestion de la sous-traitance, management de la performance SI.
    \item \textbf{Gestion de la relation avec les utilisateurs} — accompagnement au changement,
    support utilisateurs, communication, formation, gestion du niveau de service.
    \item \textbf{Développement applicatif} — projets d'informatisation, intégration, évolution
    des applications métiers.
    \item \textbf{Maintenance applicative} — suivi d'exploitation, \textbf{gestion des anomalies},
    documentation, qualité et éthique des traitements.
    \item \textbf{Gestion de l'infrastructure technique} — postes, serveurs, stockage, réseaux,
    télécoms, sécurité.
\end{enumerate}

Selon le référentiel MEF tel que présenté dans les travaux appliqués à la DCSI~\cite{meghari2014},
chaque rubrique se décline en \textbf{pratiques} puis en \textbf{activités} élémentaires (l'ordre
de grandeur usuel est de l'ordre de vingt pratiques et quatre-vingts activités pour le seul
référentiel «\,fonction SI\,», selon la version du questionnaire). À titre illustratif, sans visée
normative pour ce PFE, on retrouve notamment sous le \emph{pilotage} : stratégie et schéma
directeur, gestion de la sous-traitance, management de la performance SI, veille et audit ; sous
la \emph{relation utilisateurs} : accompagnement du changement, support utilisateurs, communication,
formation ; sous le \emph{développement applicatif} : conduite de projets d'informatisation et
intégration ; sous la \emph{maintenance applicative} : suivi d'exploitation, \textbf{gestion des
anomalies}, documentation, qualité ; sous l'\emph{infrastructure} : postes et serveurs, stockage,
réseaux et télécoms, sécurité. Cette granularité montre que la simple mise en place d'un outil de
ticketing ne «\,résout\,» qu'une fraction explicite du référentiel — mais une fraction
\textbf{structurante} pour la mesure et l'amélioration continue.

\medskip
\noindent\textbf{Lien avec le présent mémoire.} Le projet GIA ne constitue pas une application
empirique du MEF (pas d'administration des questionnaires ni de baromètres dans le cadre de ce
PFE). En revanche, le référentiel ci-dessus \textbf{justifie} la pertinence d'un outil dédié à la
\textbf{gestion des incidents applicatifs} et à la \textbf{relation utilisateurs} : sans
traçabilité ni indicateurs, le pôle «\,maintenance applicative\,» et la satisfaction des «\,clients
internes\,» des branches restent difficiles à piloter. La plateforme GIA répond précisément à ce
manque opérationnel identifié au sein du Groupe Informatique de la Branche Carburants.

\section{Enjeux contemporains et gestion des incidents}
La densification des applications, la dépendance aux réseaux et les référentiels de bonnes
pratiques (ITIL, ISO~20000, etc., cités comme cadres conceptuels) recommandent la formalisation
des incidents, la priorisation et le suivi des délais. GIA matérialise ces principes sur un
périmètre ciblé : incidents applicatifs, cycle de vie explicite et journal d'audit.

\section{Conclusion}
Ce chapitre a posé les fondations conceptuelles — définitions, typologies, acteurs, organisation
de la fonction SI — et a situé le projet dans la continuité des travaux sur Naftal et la DCSI.
Les repères MEF rappellent que le support et la gestion des anomalies sont une composante à part
entière de la performance du système d'information d'entreprise. Le chapitre suivant présente
Naftal, son environnement, la DCSI au niveau opérationnel et l'étude de l'existant ayant motivé
GIA.
